\documentclass[11pt,a4paper]{article}
\usepackage[utf8]{inputenc}
\usepackage[margin=1in]{geometry}
\usepackage{graphicx}
\usepackage{booktabs}
\usepackage{hyperref}
\usepackage{xcolor}
\usepackage{float}
\usepackage{fancyhdr}
\usepackage{titlesec}
\usepackage{listings}
\usepackage{courier}

% Header/Footer
\pagestyle{fancy}
\fancyhf{}
\rhead{Assignment 4 - GitHub Actions CI}
\lhead{MLOps Course}
\rfoot{Page \thepage}

% Colors
\definecolor{bugred}{RGB}{178, 34, 34}
\definecolor{fixgreen}{RGB}{34, 139, 34}
\definecolor{codebg}{RGB}{245, 245, 245}
\definecolor{codeframe}{RGB}{200, 200, 200}

% Listings style for YAML
\lstdefinestyle{yamlstyle}{
    backgroundcolor=\color{codebg},
    basicstyle=\ttfamily\small,
    breaklines=true,
    frame=single,
    rulecolor=\color{codeframe},
    numbers=left,
    numberstyle=\tiny\color{gray},
    xleftmargin=2em,
    framexleftmargin=1.5em,
    tabsize=2,
    showstringspaces=false,
    commentstyle=\color{gray},
    keywordstyle=\color{blue},
}

\title{\textbf{Assignment 4: GitHub Actions CI Pipeline}\\[0.5em]
\large Bug Analysis \& Corrected Workflow}
\author{Student A}
\date{March 18, 2026}

\begin{document}

\maketitle

% ============================================================
\section{Introduction}
% ============================================================

This report documents the bugs found in the provided GitHub Actions workflow file (\texttt{ml-pipeline.yml}) and the fixes applied to produce a valid, fully functional CI pipeline. The corrected pipeline runs on every push to all branches \textbf{except} \texttt{main}, installs project dependencies, runs a linter check, validates the model environment, and uploads the repository \texttt{README.md} as a GitHub Actions artifact.

% ============================================================
\section{Original (Buggy) YAML}
% ============================================================

The YAML file as originally provided is shown below. Line numbers are included for reference in the bug table.

\begin{lstlisting}[style=yamlstyle, caption={Original ml-pipeline.yml (buggy)}]
# .github/workflows/ml-pipeline.yml
name: ML Model CI

on:
push:
branches: main
pull_request:

jobs:
validate-and-test:
runs-on: ubuntu-latest
steps:
- name: Set up Python
uses: actions/setup-python@v5
with:
python-version: '3.10'

- name: Install Dependencies
run: pip install -r requirements.txt

- name: Linter Check

- name: Model Dry Test
run: |
python -c "import torch; print('Model environment ready!')"
\end{lstlisting}

% ============================================================
\section{Bugs Identified \& Fixes}
% ============================================================

\begin{table}[H]
\centering
\caption{Bug Report Summary}
\begin{tabular}{@{}c p{4.5cm} p{6.5cm}@{}}
\toprule
\textbf{\#} & \textbf{Bug Description} & \textbf{Fix Applied} \\
\midrule
\textcolor{bugred}{\textbf{1}} &
\textcolor{bugred}{\textbf{Missing \texttt{actions/checkout} step}} --- The workflow never checks out the repository, so files like \texttt{requirements.txt} are not available on the runner. &
\textcolor{fixgreen}{Added a new first step:} \texttt{uses: actions/checkout@v4} before all other steps. \\
\addlinespace
\textcolor{bugred}{\textbf{2}} &
\textcolor{bugred}{\textbf{Incorrect indentation under \texttt{on:}}} --- \texttt{push:} and \texttt{pull\_request:} are at the same level as \texttt{on:}, making them invalid trigger keys. &
\textcolor{fixgreen}{Indented \texttt{push:} and \texttt{pull\_request:}} two spaces under \texttt{on:}, and \texttt{branches-ignore:} two spaces under \texttt{push:}. \\
\addlinespace
\textcolor{bugred}{\textbf{3}} &
\textcolor{bugred}{\textbf{\texttt{branches: main} is not a list}} --- GitHub Actions expects \texttt{branches} to be a YAML list, not a bare string. &
\textcolor{fixgreen}{Changed to list syntax:} \texttt{branches-ignore:} with \texttt{- main} (also satisfies the assignment requirement to exclude \texttt{main}). \\
\addlinespace
\textcolor{bugred}{\textbf{4}} &
\textcolor{bugred}{\textbf{Linter Check step is empty}} --- The step has a \texttt{name:} but no \texttt{run:} or \texttt{uses:} key, causing workflow validation to fail. &
\textcolor{fixgreen}{Added a \texttt{run:} command:} installs \texttt{flake8} and runs it with critical error selectors (\texttt{E9,F63,F7,F82}). \\
\addlinespace
\textcolor{bugred}{\textbf{5}} &
\textcolor{bugred}{\textbf{Missing artifact upload step}} --- The assignment requires uploading \texttt{README.md} as a GitHub artifact named \texttt{project-doc}. &
\textcolor{fixgreen}{Added a final step} using \texttt{actions/upload-artifact@v4} with \texttt{name: project-doc} and \texttt{path: README.md}. \\
\bottomrule
\end{tabular}
\end{table}

% ============================================================
\section{Corrected YAML}
% ============================================================

\begin{lstlisting}[style=yamlstyle, caption={Corrected ml-pipeline.yml}]
# .github/workflows/ml-pipeline.yml
name: ML Model CI

on:
  push:
    branches-ignore:
      - main
  pull_request:

jobs:
  validate-and-test:
    runs-on: ubuntu-latest
    steps:
      - name: Checkout Repository
        uses: actions/checkout@v4

      - name: Set up Python
        uses: actions/setup-python@v5
        with:
          python-version: '3.10'

      - name: Install Dependencies
        run: pip install -r requirements.txt

      - name: Linter Check
        run: |
          pip install flake8
          flake8 . --count --select=E9,F63,F7,F82 --show-source --statistics

      - name: Model Dry Test
        run: |
          python -c "import torch; print('Model environment ready!')"

      - name: Upload Project Documentation
        uses: actions/upload-artifact@v4
        with:
          name: project-doc
          path: README.md
\end{lstlisting}

% ============================================================
\section{GitHub Actions --- Successful Run}
% ============================================================

\subsection{Actions Tab Screenshot}

\begin{figure}[H]
\centering
\fbox{\parbox{0.85\textwidth}{\centering\vspace{3cm}\textit{[Placeholder --- Insert screenshot of the GitHub Actions tab showing a successful (green) run here]}\vspace{3cm}}}
\caption{GitHub Actions tab showing a successful pipeline run}
\label{fig:actions_green}
\end{figure}

\subsection{Repository Link}

The public GitHub repository for this project is available at:

\vspace{0.5em}
\noindent\url{https://github.com/YOUR_USERNAME/YOUR_REPO_NAME}

\textit{(Replace with your actual repository URL.)}

% ============================================================
\section{Conclusion}
% ============================================================

Five bugs were identified and fixed in the original workflow file. The corrected pipeline:

\begin{itemize}
    \item Checks out the repository code before any other step.
    \item Uses correct YAML indentation and list syntax for trigger configuration.
    \item Runs on every \texttt{push} to all branches \textbf{except} \texttt{main}, and on all pull requests.
    \item Includes a working linter check using \texttt{flake8}.
    \item Uploads the project \texttt{README.md} as a GitHub artifact named \texttt{project-doc}.
\end{itemize}

\end{document}
